\section{Binomial moments}\label{sec:binom}
%
%
%

We begin by proving a special case of Theorem \ref{cycles_sets_expectation},
where each set $I_j$ is a singleton.  This is Theorem 7 in \cite{watterson}.

\begin{lem}\label{cycles}
Let $m_1,\ldots,m_n$ be non-negative integers with $m_1+2m_2+\cdots+nm_n\le n$.  Then
 \[
\E \prod_{j=1}^n \binom{C_j(\sigma)}{m_j} = \prod_{j=1}^n \frac{(1/j)^{m_j}}{m_j!}.
\]
If  $m_1+2m_2+\cdots+nm_n > n$, then the left side is zero.
\end{lem}
\begin{proof}
The second assertion is obvious, since the only way for the product on the left to be positive is for the sum of the cycle lengths to exceed $n$. Now assume that
$m_1+2m_2+\cdots+nm_n\le n$. 
 The number of ways of choosing from $[n]$ a disjoint collection 
  of $m_1$ $1-$element sets, $m_2$  $2-$element sets,
$\ldots$, $m_n$ $n-$element sets is equal to
\[
 \binom{n}{\underbrace{1 \cdots 1}_{m_1} \underbrace{2 \cdots 2}_{m_2} \cdots \underbrace{n\cdots n}_{m_n} \, t} 
 \frac{1}{m_1! \cdots m_n!} = \frac{n!/t!}{\prod_{j=1}^n (j!)^{m_j} m_j!},
 \]
where $t=n-(m_1+2m_2+\cdots+nm_n)$.
A $k$-element set may be arranged into a cycle in $(k-1)!$ ways.
Thus, 
the number of ways to arrange the elements of these sets into cycles is $(0!)^{m_1} (1!)^{m_2} \cdots (n-1)!^{m_k}$.
Finally, the $t$ elements not used in any of these cycles may be permuted in $t!$ ways.
\end{proof}
